\documentclass{standalone}
\usepackage{tikz}
\usepackage{siunitx} % Para la notación de unidades

\begin{document}

\begin{tikzpicture}
    % Núcleo
    \node[circle, draw, fill=gray!20, minimum size=2cm, inner sep=0pt] (nucleus) {};
    
    % Definir radio
    \def\R{1} % Radio del núcleo 

    % Posiciones de protones y neutrones
    \foreach \x in {-0.75, -0.45, -0.15, 0.15, 0.45, 0.75} {
        \foreach \y in {-0.75, -0.45, -0.15, 0.15, 0.45, 0.75} {
            % Verificar si el punto está dentro del núcleo
            \pgfmathsetmacro{\dist}{sqrt(\x*\x + \y*\y)}
            \ifdim\dist pt<\R pt
                \pgfmathsetmacro{\color}{mod(\x*\y*10, 3)} % Color alternante modificado
                \ifnum \color=0
                    \node[circle, draw, fill=blue, minimum size=0.2cm, inner sep=0pt] at (\x, \y) {n⁰}; % Neutrones
                \else
                    \node[circle, draw, fill=red, minimum size=0.2cm, inner sep=0pt] at (\x, \y) {p⁺}; % Protonos
                \fi
            \fi
        }
    }
    \draw[latex-latex] (-1,-1.25)--(1,-1.25);
    \node[below] at (0,-1.25){\SI{1e-14}{m}};
\end{tikzpicture}

\end{document}
